\documentclass[10pt,a4paper]{amsart}
\usepackage[margin=19mm]{geometry}
\usepackage{amsmath,amssymb,mathtools}
\usepackage[scr=rsfs,cal=euler]{mathalfa}
\usepackage{xfrac}
\usepackage[unicode,hidelinks,bookmarks=false]{hyperref}
\newcommand{\ie}{i.e.\ }
\newcommand{\cf}{cf.\ }
\newcommand{\resp}{respectively\ }
\newcommand{\Subsection}{\subsection}
\providecommand{\nobreakdash}{\nobreak\hbox{-}}
\setlength{\emergencystretch}{2em}
\multlinegap0pt
\usepackage{amssymb}
\usepackage[shortlabels]{enumitem}
\usepackage{tikz}
\newtheorem{proposition}{Proposition}
\newtheorem{theorem}{Theorem}
\DeclareMathOperator{\Aut}{Aut}
\DeclareMathOperator{\Cont}{Cont}
\title{Structures in topological recursion relations}
\author{Felix Janda, Xin Wang}
\date{Source: arXiv:2601.20673v1 (CC BY 4.0)}
\begin{document}
\maketitle

\noindent\emph{Source and licence.} Structures in topological recursion relations. Version arXiv:2601.20673v1 (CC BY 4.0), distributed by arXiv under the Creative Commons Attribution 4.0 International licence, http://creativecommons.org/licenses/by/4.0/, as recorded in the arXiv licence metadata for this exact version and shown on the abstract page https://arxiv.org/abs/2601.20673v1. Full licence notice: \texttt{licenses/arxiv-2601.20673-CC-BY-4.0.txt}.\par\smallskip

\noindent\textit{Editorial scope.} Counted unit: the article's theorem thm:algo-inter (source0.tex lines 1682--1689). The article's complete original proof of it is delivered with it (source0.tex lines 1690--1781, one \texttt{proof} environment of 92 lines). No mathematics is written, repaired or re-derived: the statement and the proof are the article's own lines, in the article's own notation, and no retained line is substituted or re-flowed.  Retained uncounted as the source-local context the proof needs: lines 636--648; lines 1633--1635.  Nothing was omitted from any retained span, so the package carries no omission record.  No retained line contains a citation, so the wrapper carries no bibliography and no bibliography entry.  No retained line contains a figure environment, an image-inclusion command or drawing code, so no image is delivered and the package declares no asset dependency.  Editorial formatting only, in addition: an AMS \texttt{amsart} preamble that restates the article's own declarations and macros used by the retained lines, namely the environments \texttt{proposition}, \texttt{theorem} on separate counters, together with the standard packages those lines need.  The retained environments print in this document as Proposition 1, Theorem 1 in order of appearance, whereas the article prints the same items with the numbering of its own full document.  The complete native source of the article as distributed in the arXiv e-print for this exact version is bundled unchanged as \texttt{sources/193541/source0.tex}.
\begin{proposition}\label{prop:omega_g,m-formula}
  For any monomial $M$ of degree $D \le 2g + 1$, we have
  \begin{align}\label{eqn:omega_g,m-formula}
    \Omega_{g, M}
    = \sum_{\Gamma\in G_{g,n+1}} \sum_{\substack{m\colon V \to \mathbb Z_{\ge 0} \\ \sum_v m(v) = 2g + 1 - D}}
    \frac{(m(v^*) + 1)(2g + 1 - D)!}{|\Aut(\Gamma)|}
    \prod_{v \in V} \frac{(2g(v) - 3 + n(v) + m(v))!}{(2g(v) - 3 + n(v))!} \cdot C_{\Gamma, M, m},
  \end{align}
  where $v^*$ is the vertex containing the $(n + 1)$st leg, and where
  $C_{\Gamma, M, m}$ is the coefficient of
  $M \cdot x_{v^*} \prod_v x_v^{m(v)}$ of the degree-$(g + 1)$ part of
  the class $\mathcal{D}_{\Gamma}(a_2, \dotsc, a_n, \{x_v\})$.
\end{proposition}

\begin{align}\label{eqn:int-trr-3g-3+n}
\int_{\overline{\mathcal{M}}_{g,n}}\widehat{\Omega}_{g,M} =0. 
\end{align}

\begin{theorem}\label{thm:algo-inter}
  For any $\sum_{i=1}^{n}k_i=3g-3+n$, we have
\begin{align}\label{eqn:deg-3g-3+n-push-algo-int}
\int_{\overline{\mathcal{M}}_{g,n}}  \prod_{i=1}^{n}{\psi_i}^{k_i}   =\sum_{\substack{\Gamma\in G_{g,n}\\ |E(\Gamma)|\geq1}}\sum_{\{\alpha_{v}:v\in V(\Gamma)\}} c_{(\Gamma,\{\alpha_v\}_{v\in V(\Gamma)})}\cdot \prod_{v\in V(\Gamma)}\int_{\overline{\mathcal{M}}_{g(v),n(v)}}\alpha_{v}  
\end{align}
where  $\alpha_v$ ranges over all possible  monomials of tautological $\psi$ classes on $\overline{\mathcal{M}}_{g(v),n(v)}$.
Each coefficient $c_{(\Gamma,\{\alpha_v\}_{v\in V(\Gamma)})}$ is a rational number coming from the weight on stable graphs in Pixton's double ramifcation formula and factors from string equation and dilaton equation.        
\end{theorem}

\begin{proof}
The proof is similar to the proof of Proposition \ref{prop:omega_g,m-formula}.   First, the multiplication by $\psi_{n+1} \cdot \dots \cdot \psi_N$
  in the definition of $\widehat{\Omega}_{g, M}$ ensures that we only need to
  consider $N$-pointed stable graphs obtained by distributing
  $N-n = 2(3g-3+n)-D$ legs to an genus-$g$, $n$-pointed stable graph
  $\Gamma$.
  Let $m(v)$ be the number of extra legs at a vertex $v$ of $\Gamma$.
  Then, there are
  \begin{equation*}
    \frac{(2(3g-3+n)-D)!}{\prod_v m(v)!}
  \end{equation*}
  many such distributions.
  Because of the identity
  \begin{equation*}
    [p(x_1 + \dotsb + x_m)]_{x_1 \cdot \dotsb \cdot x_m}
    = m! [p(x)]_{x^m}
  \end{equation*}
  for any polynomial $p(x)$, we may set $x_v$ for $v$ be the
  sum of the variables $a_i$ for $i$ corresponding to extra markings
  at $v$.
  This leads to the additional factors
  \begin{equation*}
\prod_{v} m(v)!.  \end{equation*}
  Finally, applying the dilaton equation repeatedly to compute the
  push-forward under the map forgetting the last $2(3g-3+n)-D$ markings
  leads to the factor
  \begin{equation*}
    \frac{(2g(v) - 3 + n(v) + m(v))!}{(2g(v) - 3 + n(v))!}
  \end{equation*}
  for every vertex $v$.
  Putting all of these together yields the formula
  \begin{align}\label{eqn:hat-omega_g,m-formula}
   \widehat{\Omega}_{g, M}
    = \sum_{\Gamma\in G_{g,n}} \sum_{\substack{m\colon V \to \mathbb Z_{\ge 0} \\ \sum_v m(v) = 2(3g-3+n)-D}}
    \frac{(2(3g-3+n)-D)!}{|\Aut(\Gamma)|}
    \prod_{v \in V} \frac{(2g(v) - 3 + n(v) + m(v))!}{(2g(v) - 3 + n(v))!} \cdot \widehat{C}_{\Gamma, M, m},
  \end{align}
  where 
  $\widehat{C}_{\Gamma, M, m}$ is the coefficient of
  $M \cdot \prod_v x_v^{m(v)}$ of the degree $3g-3+n$ part of
  the class $\widehat{\mathcal{D}}_{\Gamma}(a_2, \dotsc, a_n, \{x_v\})$.


Combining equations~\eqref{eqn:int-trr-3g-3+n} and \eqref{eqn:hat-omega_g,m-formula}, we obtain
\begin{align}\label{eqn:graph-sum-deg-3g-3+n-int}
\sum_{\Gamma\in G_{g,n}} \sum_{\substack{m\colon V \to \mathbb Z_{\ge 0} \\ \sum_v m(v) = 2(3g-3+n)-D}}
    \frac{(2(3g-3+n)-D)!}{|\Aut(\Gamma)|}
    \prod_{v \in V} \frac{(2g(v) - 3 + n(v) + m(v))!}{(2g(v) - 3 + n(v))!} \cdot \int_{\overline{\mathcal{M}}_{g,n}}\widehat{C}_{\Gamma, M, m}=0.
  \end{align}
  Finally, the above equation~\eqref{eqn:graph-sum-deg-3g-3+n-int} leads to equation~\eqref{eqn:deg-3g-3+n-push-algo-int}: letting $M=\prod_{i=2}^{n}a_i^{2k_i}$,
  the contribution of the trivial stable graph $\Gamma_{g,[n]}$ in equation~\eqref{eqn:graph-sum-deg-3g-3+n-int} becomes
  \begin{align}\label{eqn:contr-trivial-inter-number}
  &(6g-6+2n-D)!\frac{(2g-3+n+2(3g-3+n)-D)!}{(2g-3+n)!} 
\nonumber\\&\hspace{10pt}\cdot\int_{\overline{\mathcal{M}}_{g,n}}\left[\widehat{\mathcal{D}}_{\Gamma_{g,[n]}}^{3g-3+n}(a_2,\dots,a_n,\{x\})\right]_{\prod_{i=2}^{n}a_i^{2k_i}\cdot x^{6g-6+2n-D}}
\nonumber\\=&(6g-6+2n-D)!\frac{(8g-9+3n-D)!}{(2g-3+n)!} 
\int_{\overline{\mathcal{M}}_{g,n}}\left[e^{\frac{1}{2}(x+\sum_{i=2}^{n}a_i)^2\psi_1}\prod_{i=2}^{n}e^{\frac{1}{2}a_i^2\psi_i}\right]_{\prod_{i=2}^{n}a_i^{2k_i}\cdot x^{6g-6+2n-D}}
\nonumber\\=&(6g-6+2n-D)!\frac{(8g-9+3n-D)!}{(2g-3+n)!} 
\nonumber\\&\cdot
\sum_{\sum_{i=1}^{n}l_i=3g-3+n}\prod_{i=1}^{n}\frac{1}{2^{l_i}l_i!}
\int_{\overline{\mathcal{M}}_{g,n}}\prod_{i=1}^{n}\psi_i^{l_i}\left[(x+\sum_{i=2}^{n}a_i)^{2l_1}\right]_{\prod_{i=2}^{n}a_i^{2k_i-2l_i}\cdot x^{6g-6+2n-D}}
\nonumber\\=&(6g-6+2n-D)!\frac{(8g-9+3n-D)!}{(2g-3+n)!} 
\nonumber\\&\cdot
\sum_{\sum_{i=1}^{n}l_i=3g-3+n}\prod_{i=1}^{n}\frac{1}{2^{l_i}l_i!}
\int_{\overline{\mathcal{M}}_{g,n}}\prod_{i=1}^{n}\psi_i^{l_i}
\binom{2l_1}{2k_2-2l_2,\dots,2k_n-2l_n,6g-6+2n-2\sum_{i=2}^{n}k_i}.
  \end{align}
The contribution of all stable graphs with at least one edge in equation~\eqref{eqn:graph-sum-deg-3g-3+n-int} is equal to 
\begin{align}\label{eqn:contr-other-graph-inter-number}
&\sum_{\Gamma\in G_{g,n}: |E(\Gamma)|\geq1}\sum_{\substack{m:V\rightarrow\mathbb{Z}_{\geq0}\\
\sum_{v}m(v)=2(3g-3+n)-D}}\frac{(6g-6+2n-D)!}{|\Aut(\Gamma)|}
\nonumber\\&\hspace{20pt}\cdot\prod_{v\in V(\Gamma)}\frac{(2g(v)-3+n(v)+m(v))!}{(2g(v)-3+n(v))!}\int_{\prod_{v}\overline{\mathcal{M}}_{g(v),n(v)}}\left[\widehat{\mathcal{D}}_{\Gamma}^{3g-3+n}(a_2,\dots,a_n,\{x_v\})\right]_{\prod_{i=2}^{n}a_i^{2k_i}\cdot\prod_{v\in V(\Gamma)}x_v^{m_v}}.    \end{align}
Plugging equations~\eqref{eqn:contr-trivial-inter-number} and \eqref{eqn:contr-other-graph-inter-number} into equation~\eqref{eqn:graph-sum-deg-3g-3+n-int}, we obtain
\begin{align*}
\int_{\overline{\mathcal{M}}_{g,n}}\prod_{i=1}^{n}\psi_i^{k_i} +\sum_{(l_2,\dots,l_n)<(k_2,\dots,k_n)}c'_{l_2,l_3,\dots,l_n}
\int_{\overline{\mathcal{M}}_{g,n}}\prod_{i=1}^{n}\psi_i^{l_i} =\text{boundary terms}
\end{align*}
where $(l_2,\dots,l_n)<(k_2,\dots,k_n)$ means the monomials of $\psi$ classes $\prod_{j=1}^{n}\psi_j^{l_j}$ strictly lower than  $\prod_{j=1}^{n}\psi_j^{k_j}$. Here all $c'_{k_2,k_3,\dots,k_n}$ are rational numbers determined from equation~\eqref{eqn:contr-trivial-inter-number}.
Finally, by induction on $(k_2,\dots,k_n)$, we obtain equation~\eqref{eqn:deg-3g-3+n-push-algo-int}. 
% so\dots
% \begin{align*}
% 0= &\int_{\overline{\mathcal{M}}_{g,N}}\left[\mathcal{D}_{g,N}^{3g-3+n}(a_2,\dots,a_N)\right]_{a_2^{2k_2}\dots a_n^{2k_n}\cdot a_{n+1}\dots a_N}  \cdot \psi_{n+1}\dots\psi_{N}
% \\=&\int_{\overline{\mathcal{M}}_{g,N}}\left[\prod_{i=1}^{N}e^{\frac{1}{2}a_i^2\psi_i}\right]_{a_2^{2k_2}\dots a_n^{2k_n}\cdot a_{n+1}\dots a_N}\cdot \psi_{n+1}\dots\psi_{N}+\text{boundary terms}
% \\=&\frac{(2g-3+N)!}{(2g-3+n)!}\int_{\overline{\mathcal{M}}_{g,n}}\left[\prod_{i=1}^{n}e^{\frac{1}{2}a_i^2\psi_i}\right]_{a_2^{2k_2}\dots a_n^{2k_n}\cdot a_{n+1}\dots a_N}
% \\&+\sum_{\Gamma\in G_{g,N}: |E(\Gamma)|\geq1}\int_{\overline{\mathcal{M}}_{g,N}}\Cont_{\Gamma}\left[\mathcal{D}_{g,N}^{3g-3+n}(a_2,\dots,a_N)\right]_{a_2^{2k_2}\dots a_n^{2k_n}\cdot a_{n+1}\dots a_N}\psi_{n+1}\dots\psi_{N} \\=&\frac{(2g-3+N)!}{(2g-3+n)!}\int_{\overline{\mathcal{M}}_{g,n}}\left[\prod_{i=1}^{n}e^{\frac{1}{2}a_i^2\psi_i}\right]_{a_2^{2k_2}\dots a_n^{2k_n}\cdot a_{n+1}\dots a_N}
% \\&+\sum_{\Gamma\in G_{g,n}: |E(\Gamma)|\geq1}\sum_{\substack{m\colon V\rightarrow\mathbb{Z}_{\geq0}\\
% \sum_{v}m(v)=2(3g-3+n)-D}}\frac{(6g-6+2n-D)!}{|\Aut(\Gamma)|}
% \\&\hspace{30pt}\cdot\prod_{v\in V(\Gamma)}\int_{\overline{\mathcal{M}}_{g(v),n(v)}}\left[\mathcal{D}_{\Gamma}^{3g-3+n}(a_2,\dots,a_n,\{x_v\})\right]_{a_2^{2k_2}\dots a_n^{2k_n}\cdot\prod_{v\in V(\Gamma)x_v^{m_v}}}\prod_{v}\frac{(2g(v)-3+n(v)+m(v))!}{(2g(v)-3+n(v))!}
% \\=&\frac{(2g-3+N)!}{(2g-3+n)!}
% \int_{\overline{M}_{g,n}}\psi_{1}^{3g-3+n-k_2-\dots-k_n}\psi_2^{k_2}\dots\psi_{n}^{k_n}+\dots
% .
% \end{align*}    
\end{proof}

\end{document}
